\documentclass[11pt]{article}

\usepackage[margin=1in]{geometry}
\usepackage{amsmath}
\usepackage{amssymb}
\usepackage{booktabs}
\usepackage[hidelinks]{hyperref}

\title{Distributed Bundle Adjustment in \texttt{client\_acc.py}}
\author{}
\date{}

\begin{document}
\maketitle

\section{Overview}

The implementation coordinates bundle-adjustment subproblems over several
landmark clusters.  Each cluster owns a disjoint set of landmarks, but cameras
may occur in several clusters.  Consequently, optimized landmarks can be
written directly back to their global rows, whereas the duplicate camera
estimates must be reconciled by a consensus step.

The outer algorithm is a variable-metric Douglas--Rachford splitting (DRS)
iteration with a Nesterov-like extrapolation of its cluster reference states.
The local nonlinear least-squares problems are solved by the C++ workers; the
Python coordinator dispatches those solves, merges their results, tests
accelerated and unaccelerated outer steps, and maintains recovery state.

\section{Notation and Code Variables}

For cluster $i$, let
\begin{align*}
u_i &\quad\text{denote its optimized local camera copy},\\
s_i &\quad\text{denote its proximal reference state},\\
U_i &\quad\text{denote its block-diagonal camera metric},\\
\ell_i &\quad\text{denote its owned landmarks},
\end{align*}
and let $v$ denote the global consensus camera vector.  Every camera has nine
parameters, so each nonzero camera block of $U_i$ is a $9\times9$ matrix.

\begin{center}
\begin{tabular}{ll}
\toprule
Mathematical object & Python representation \\
\midrule
$u_i$ & \verb|poses_in_cluster[i]| \\
$s_i$ & \verb|poses_s_in_cluster[i]| \\
$s_i^{+}$ & \verb|poses_s_in_cluster_pre[i]| \\
$v$ & \verb|poses_v| \\
$U_i$ & \verb|Ul_in_cluster[i]| or \verb|Up_cluster[i]| \\
$\ell_i$ & the rows owned by cluster $i$ in \verb|landmarks| \\
\bottomrule
\end{tabular}
\end{center}

Only cameras present in a cluster are active in that cluster's vectors and
metric.  The Python arrays retain global camera indexing, while absent rows
are masked to zero.

\section{Local Cluster Problems}

The function \verb|prox_f_push_pull| sends one request to every worker and
collects the replies in arbitrary arrival order.  Conceptually, cluster $i$
solves the proximal bundle-adjustment problem
\begin{equation}
  (u_i,\ell_i)
  \approx
  \underset{u,\ell}{\operatorname{argmin}}
  \left\{
    f_i(u,\ell)
    + \frac{1}{2}(u-s_i)^{\mathsf T}U_i(u-s_i)
  \right\},
  \label{eq:local-prox}
\end{equation}
where $f_i$ is the reprojection objective for the observations assigned to
cluster $i$.  The C++ server implements the second term as a camera prior and
returns the optimized cameras, optimized landmarks, objective value, and one
$9\times9$ metric block for every local camera.

Replies carry a cluster identifier as well as run and phase identifiers.  The
coordinator therefore associates out-of-order replies with the correct
cluster and rejects stale replies left by an earlier asynchronous batch.

\section{How Cluster Results Are Merged}

\subsection{Landmarks}

Clustering is performed by landmark: all observations of a landmark are
assigned to one cluster.  A landmark consequently has a single optimizing
owner and requires no averaging.  On receipt of a cluster reply,
\verb|prox_f_push_pull| writes the returned landmark values directly into the
corresponding global rows of \verb|landmarks|.

\subsection{Cameras}

A camera may observe landmarks in several clusters and can therefore have
several local estimates.  The function \verb|average_cameras_new| embeds each
cluster metric into the global camera space and forms the reflected DRS state
\begin{equation}
  r_i = 2u_i-s_i.
\end{equation}
It then computes the block-metric consensus
\begin{equation}
  v =
  \left(\sum_i U_i\right)^{-1}
  \sum_i U_i(2u_i-s_i).
  \label{eq:consensus}
\end{equation}
Because all $U_i$ are block diagonal by camera, Equation~\eqref{eq:consensus}
is evaluated independently for each camera.  For camera $c$ it becomes
\begin{equation}
  v_c =
  \left(\sum_{i:\,c\in i}U_{i,c}\right)^{-1}
  \sum_{i:\,c\in i}U_{i,c}(2u_{i,c}-s_{i,c}).
\end{equation}
This is not an arithmetic average.  A cluster contributes through a full
$9\times9$ matrix, allowing its estimate to have different weights in
different camera-parameter directions.

\section{Plain Douglas--Rachford Update}

After obtaining $v$, the unaccelerated update of each reference state is
\begin{equation}
  s_i^{+}=s_i+\tau(v-u_i).
  \label{eq:drs-update}
\end{equation}
The code currently sets $\tau=1$, giving $s_i^{+}=s_i+v-u_i$.  The proposed
plain state is stored in \verb|poses_s_in_cluster_pre[i]|.  If all local
camera copies agree with the consensus, then $v-u_i=0$ and the reference
state no longer moves.

\section{Nesterov-Like Acceleration}

\subsection{Accelerated State}

Acceleration is applied to the outer DRS reference states, not directly to
the landmarks, consensus cameras, or internal Ceres iterations.  The code
stacks all cluster states into one vector,
\begin{equation}
  s_k = [s_{1,k}^{\mathsf T},\ldots,s_{K,k}^{\mathsf T}]^{\mathsf T},
\end{equation}
whose allocated size is $9K$ times the total number of cameras.  Let
\begin{equation}
  p_k=s_k^{+}-s_k
\end{equation}
be the plain DRS step.  The active momentum recurrence is
\begin{align}
  \beta_k &= \frac{k-r-1}{k-r+2},
  \label{eq:beta}\\
  d_k &= p_k+\beta_k d_{k-1},
  \label{eq:direction}\\
  s_k^{\mathrm{acc}} &= s_k+d_k,
  \label{eq:accelerated-state}
\end{align}
where $r$ is \verb|resetIt|, the iteration index of the latest restart.  Thus
the accelerated trial can equivalently be written as
\begin{equation}
  s_k^{\mathrm{acc}}=s_k^{+}+\beta_k d_{k-1}.
\end{equation}

The source also contains a conventional FISTA-style recurrence using
\verb|lambda_0|, \verb|lambda_1|, and \verb|gamma|, but it is inside an
\verb|if False| block and is not active.  Therefore, Equations
\eqref{eq:beta}--\eqref{eq:accelerated-state}, rather than the disabled FISTA
formula, describe the acceleration actually used by the coordinator.

\subsection{Accelerated Trial and Plain Fallback}

After the first iterations following a restart, the code normally evaluates
two proximal centers.  It constructs
\begin{equation}
  \widehat{s}_k(t)
  =t s_k^{+}+(1-t)s_k^{\mathrm{acc}}.
  \label{eq:trial-interpolation}
\end{equation}
With two line-search iterations, the first trial uses $t=0$ and is therefore
the accelerated state.  If that trial is not accepted, the second uses $t=1$
and is the plain DRS state.  For every trial the coordinator performs a full
set of local proximal solves, recomputes the metric consensus, and evaluates
the resulting merit values.

Variables carrying the suffix \verb|_bfgs| hold these trial states and
results.  The active code does not perform a BFGS inverse-Hessian update, so
the suffix is historical and should not be interpreted as the optimization
method used here.

\section{Merit Function, Acceptance, and Recovery}

The function \verb|cost_DRE| computes the metric Douglas--Rachford envelope
contribution
\begin{equation}
  \mathcal{E}_i
  =\frac{1}{2}(u_i-v)^{\mathsf T}U_i
    \bigl((u_i-v)+2(u_i-s_i)\bigr).
  \label{eq:dre}
\end{equation}
The trial merit adds the local reprojection costs to the sum of these
contributions.  The implementation also evaluates the reprojection objective
at the consensus cameras, $f(v)$, and uses the maximum of this primal value
and the computed envelope merit as a safeguard.  Diagnostics include the
local cost $f(u)$ and the disagreement gap $f(v)-f(u)$.

If the accelerated trial improves the merit sufficiently, it is accepted and
the plain fallback is skipped.  Otherwise, the workers restore the landmarks
from before the accelerated trial and solve again with the plain DRS center.
The plain trial is accepted under a more permissive fallback condition unless
both the envelope and primal consensus costs trigger the rejection logic.

Repeated accelerated-trial failures increment
\verb|failedNesterovAcceleration|.  After
\verb|maxFailedNesterovAcceleration| failures, the implementation sets the
previous direction to zero and records the current iteration in
\verb|resetIt|.  This restarts both the momentum history and the coefficient
schedule in Equation~\eqref{eq:beta}.  The first two iterations after a
restart use only the plain DRS trial.

A severe rejection additionally restores the best-known consensus cameras
and landmarks, sets all local camera and reference states to that consensus,
increases the block regularization parameter, and clears the momentum.  These
operations are stability safeguards around the core DRS iteration rather
than part of the basic Nesterov formula.

\section{Algorithm Summary}

One outer iteration can be summarized as follows:
\begin{enumerate}
  \item Form the plain DRS reference state $s_k^{+}$ using
        Equation~\eqref{eq:drs-update}.
  \item Build the Nesterov-like direction $d_k$ using
        Equations~\eqref{eq:beta} and \eqref{eq:direction}.
  \item Solve all cluster proximal problems at the accelerated center
        $s_k^{\mathrm{acc}}$.
  \item Merge duplicate cameras using the block-metric consensus in
        Equation~\eqref{eq:consensus}; copy owned landmarks directly.
  \item Evaluate the DRS merit and the primal objective at the consensus.
  \item Accept the accelerated result if suitable; otherwise repeat the
        cluster solves at the plain center $s_k^{+}$.
  \item On repeated or severe failures, clear acceleration history and,
        when necessary, restore the best-known primal state.
\end{enumerate}

The essential distinction is that local bundle-adjustment solutions are not
merged as complete models.  Landmarks have unique owners, while shared camera
copies are combined through a variable-metric DRS projection.  Nesterov-like
momentum accelerates the sequence of DRS reference states and is protected by
a plain-step fallback and explicit restart logic.

\end{document}